\section{Descent from one marked spherical cut}

Let $T=4\pi^2$. The source is an oriented smooth closed homotopy
four-sphere. We use the Hilbert--Schmidt norm of the Weyl endomorphism
on the self-dual two-forms. A chosen smooth ALE compactification
$\widehat X$ has a marked normal ball at its added point. The comparison
must embed the whole protected core and the two collar buffers into a
single source. Several of the lemmas below apply to general measure
spaces or manifolds; their geometric applications use these choices.

\begin{lemma}[Strict exterior margin]
\label{lem:strict-exterior-margin}
\lean{FirstOrderTwistor.exteriorEnergyLtThreshold, FirstOrderTwistor.exteriorEnergyLtThresholdOfStrictError}
\leanok
Let $T,\delta>0$, $e\le2T-\delta$, $v\le e-x+\eta$ and
$0\le\eta\le\delta$. If $x>T$, then $v<T$. If $x\ge T$, the
same conclusion holds when $\eta<\delta$.
\end{lemma}
\begin{proof}
\leanok
At $x=T$, $e=2T-\delta$ and $\eta=\delta$ the exterior bound is $T$.
Either strict inequality in the statement lowers this upper bound.
\end{proof}

\begin{lemma}[Protected core and physical exterior]
\label{lem:cut-core-complement}
\lean{FirstOrderTwistor.FixedTruncationComparison.tendsto_complMass, FirstOrderTwistor.eventually_cappedExteriorEnergy_le_of_coreComparisons}
\leanok
Assume fixed-core image-mass comparison, a finite model of energy $x$
exhausted almost everywhere by increasing protected cores, and total
source energy tending to $e$. Suppose the exterior costs at most the
physical complement of the protected core plus a splice error tending
to zero in the iterated order. For every $\varepsilon>0$ and all
sufficiently large fixed truncations $L$, the exterior energy is at most
$e-x+\varepsilon$ for all sufficiently late source indices at $L$.
\end{lemma}
\begin{proof}
\leanok
Subtract the mass of the measurable embedded image from the total
source mass. At fixed $L$ this difference tends to $e$ minus the model
core mass. Use exhaustion and split the allowance between the remaining
model tail and the splice.
\end{proof}

\begin{lemma}[Cap form and splice estimates]
\label{lem:cut-cap-form}
\lean{FirstOrderTwistor.capPartitionMass, FirstOrderTwistor.capCutoffLoss, FirstOrderTwistor.capYamabeInfimumLowerBound, FirstOrderTwistor.smoothCapSpliceWeylBound, FirstOrderTwistor.smoothCapSpliceBachBound, FirstOrderTwistor.firstOrderCapSpliceWeylBound}
\leanok
For nonzero $L^4$ tests, let $\chi^2+\xi^2=1$ and suppose the two
localized quadratic forms have nonnegative lower coefficients $a,b$.
Supply the IMS identity with error $6\int Fu^2$ where $F\ge0$ and
$6\|F\|_2\le L$. Under the stated measurability and integrability
conditions the infimum of the tested Rayleigh quotients is at least
$\min(a,b)-L$. On a splice of volume at most $V$, a pointwise Weyl
bound $C\eta_2$ costs at most $V C^2\eta_2^2$ in squared norm.
An independent pointwise Bach bound $C'\eta_4$ costs at most
$VC'\eta_4$ in integrated norm. In the first-order model, put
$q,A$ for the old and new squared norms, with the quotient normalization.
If the normalized weighted $L^2$ splice error is at most $\Delta$ and
the cutoff has modulus at most one, the net added energy is at most
$2\sqrt q\,\Delta+\Delta^2$.
\end{lemma}
\begin{proof}
\leanok
Apply the triangle inequality to $u^2=(\chi u)^2+(\xi u)^2$ in
$L^2$ and use Cauchy--Schwarz on the cutoff loss. Divide the resulting
form bound by $\|u\|_4^2$. Integrate the two pointwise bounds on the
fixed splice. For the first-order estimate apply the $L^2$ triangle
inequality on the pair of splices and square.
\end{proof}

\begin{lemma}[One buffer, radius and index for both tests]
\label{lem:cut-ordered-selection}
\uses{lem:strict-exterior-margin,lem:cut-cap-form}
\lean{FirstOrderTwistor.selectFixedSphericalCut}
\leanok
Suppose $e\le2T-\delta$, $x\ge T$, and $0<4\varepsilon<\delta$.
Supply a cutoff loss tending to zero with width, compact core exhaustion
and collar closeness as the radius tends to zero for each width, and
source comparisons at each chosen fixed radius. Assume the same exterior
metric satisfies the cap energy and Yamabe estimates in the normal chart.
One can choose a width $\ell\ge1$, then an arbitrarily small fixed
$s>0$, so that eventually in the source index
\[
 E_{\rm ext}<e-x+4\varepsilon,\qquad E_{\rm ext}<T,\qquad
 Y_{\rm ext}>\tfrac12\min(y_0,Y_{\rm rd}/2)>0.
\]
\end{lemma}
\begin{proof}
\leanok
Choose $\ell$ to control the cutoff loss. In that width make the model
tail and collar error small at a fixed radius. Intersect the eventual
conditions for total energy, protected core comparison and splice error
at that radius. Their four allowances give the displayed energy bound.
\end{proof}

\begin{lemma}[The normal ball and the signed slice]
\label{lem:cut-normal-mark}
\lean{FirstOrderTwistor.normalCutPolarDiffeomorph, FirstOrderTwistor.chartedNormalCap_boundary, FirstOrderTwistor.chartedNormalCut_sides, FirstOrderTwistor.exists_normalCutRadius_for_buffer, FirstOrderTwistor.signedCollar_inside_capNeighborhood}
\leanok
In normal coordinates put $P_s(t,u)=se^{-t}Ju$ for $s>0$, $u\in S^3$
and an orthogonal $J$. The signed cylinder maps diffeomorphically onto punctured four-space,
with inverse given by the logarithmic radius and the normalized direction. The boundary of the marked ball $v\mapsto sJv$
is precisely the middle slice. The positive side lies inside the ball
and the negative side lies outside the closed ball. A radius as small
as desired places the ball and any fixed finite buffer in a chart
about zero. An open neighborhood of an entire middle sphere contains
a two-sided strip of positive uniform width.
\end{lemma}
\begin{proof}
\leanok
The norm is $se^{-t}$, giving signed time $-\log(\|y\|/s)$ and the
three radius tests. The largest radius on a buffer of half-length
$L$ is $se^L$. For the final assertion apply the tube lemma to
$\{0\}\times S^3$ in the inverse image of the open neighborhood.
\end{proof}

\begin{proposition}[Smooth reassembly on a marked collar]
\label{prop:cut-marked-reassembly}
\lean{FirstOrderTwistor.exists_markedExteriorRecognition, FirstOrderTwistor.exists_recognition_eqOn_capNeighborhood, FirstOrderTwistor.originalTypeOfRecognizedExterior, FirstOrderTwistor.SmoothCutGluing.toDiffeomorph_preservesOrientation}
\leanok
Given an oriented recognition $f:\widehat V\to S^4$ and an oriented
ambient extension of its recorded cap germ, a corrected recognition
agrees with the prescribed chart near the closed cap and maps its closed
exterior onto the opposite hemisphere. Supply the retained-side and
normal-ball maps, the core inverse, and the two open covers with their
common signed collar charts. When these charts satisfy the recorded
patch formulas and preserve tangent orientations on their domains, they
induce an oriented smooth diffeomorphism $\Phi:M\to\widehat X$
with those core and exterior restrictions.
\end{proposition}
\begin{proof}
\leanok
Compose $f$ with the inverse ambient extension. Compose the corrected
map with the retained-side and normal-ball charts. The resulting
exterior map and the core inverse coincide on the source overlap, and
the common collar maps onto the target overlap. Glue both maps and
their inverses on these opens. Their differentials on the patch domains
carry the given orientations.
\end{proof}

\begin{proposition}[Marked competitor and the variational residual]
\label{thm:one-cut-same-type-descent}
\uses{prop:cut-marked-reassembly,lem:scalar-component-budgets}
\lean{FirstOrderTwistor.classYamabeFour_pullback, FirstOrderTwistor.chiralFrameEnergy_pullback, FirstOrderTwistor.MarkedMetricTransfer.competitor, FirstOrderTwistor.lowerEnergyOnOriginalType, FirstOrderTwistor.excludesExactMinimum, FirstOrderTwistor.uniformDescentOnOriginalType}
\leanok
Use the specified positive-block coefficients for the energy functional
on all source metrics. On the compactified metric $g$ of energy $x$,
suppose $Y([g])>0$ and the oriented patch maps identify the two blocks
up to almost-everywhere orthogonal conjugacy. The actual pullback
$\Phi^*g$ has the same positive Yamabe constant and energy. The infimum
on $M$ for this coefficient family, denoted by $I_W(M)$, therefore
satisfies $I_W(M)\le x$. The strict scalar residual gives $I_W(M)\le x<e$ under its Ricci and
subround Yamabe hypotheses. With a supplied band-uniform deficit the
stronger inequality is $I_W(M)\le x\le e-\kappa_\delta$.
\end{proposition}
\begin{proof}
\leanok
Pullback identifies the two sets of normalized total scalar-curvature
values as the smooth logarithmic conformal factor varies. Conjugacy
preserves the squared Hilbert--Schmidt norm, and Riemannian volume
changes variables by the same diffeomorphism. Insert the resulting
metric into the source infimum and apply the strict or uniform residual.
\end{proof}

\begin{theorem}[Curvature and descent on the original smooth type]
\label{thm:cut-curvature-descent}
\uses{prop:cut-marked-reassembly,thm:one-cut-same-type-descent}
\lean{FirstOrderTwistor.geometricChiralBlock, FirstOrderTwistor.geometricChiralBlock_pullback, FirstOrderTwistor.geometricPullbackCompetitor, FirstOrderTwistor.geometricPositiveEnergies_eq_of_diffeomorph, FirstOrderTwistor.geometricDescentOnOriginalType, FirstOrderTwistor.oneCutDescent}
\leanok
Compute the positive block of the actual lowered Levi--Civita curvature
in an oriented orthonormal basis by restricting to the normalized
self-dual wedge basis and removing the scalar trace. Let $\mathcal E(M)$
be the energies so obtained from positive-Yamabe metrics and pointwise
oriented orthonormal frames with integrable squared block norm.
An oriented diffeomorphism identifies the two variational sets.
For the single-cut implication, assume $e\le 2T-\delta$, $x>T$ and
$0\le\eta\le\delta$ with $\delta>0$. On the specified exterior supply a
positive Yamabe metric of integrable chiral energy at most $e-x+\eta$.
Assume its strict recognition principle, the hemisphere germ-extension
theorem and the marked local patches for this cut. The map returned by reassembly agrees with the fixed core inverse
and pulls the compactified metric back to the given source. On the
marked cut, a positive-Yamabe compactification with integrable curvature
energy $x$ yields $x\in\mathcal E(M)$. The Route 02 strict residual
implies $\inf\mathcal E(M)\le x<e$. The uniform residual gives
$\inf\mathcal E(M)\le x\le e-\kappa_\delta$ with its extra Yamabe deficit.
\end{theorem}
\begin{proof}
\leanok
Transport the chosen tangent basis by the inverse differential. It
remains oriented and orthonormal for the pulled-back metric. Riemann
naturality identifies all four curvature slots, so the two positive
blocks have identical entries in these paired frames. Pullback preserves
integrability, the energy integral and the Yamabe value. Apply the
construction to the inverse map for the reverse inclusion of the
variational sets. Use the source infimum and the two scalar residuals.
\end{proof}

\begin{openproblem}[Realizing the marked cut from the source hypotheses]
\label{problem:cut-geometric-realization}
\notready
Produce the required smooth compactification, marked full-core embedding,
cap metric estimates and recognized exterior under the hypotheses of the
single-cut theorem. Identifying the formal infimum $\inf\mathcal E(M)$
with the intrinsic $I^+(M)$ in the manuscript requires invariance under
arbitrary changes of oriented orthonormal frame and integrability of the
intrinsic density for each smooth metric. Formal recognition still needs the topology of the
actual capped exterior and the strict sphere theorem. The expected strict
residual depends on the geometric Euler, Ricci and Yamabe inputs. General
existence of a suitable cut and further descent from a positive closed
critical class are separate questions.
\end{openproblem}
